\section{Equipo y ética}

\subsection{Gobernanza institucional del proyecto}

El proyecto se organiza bajo una estructura de gobernanza transparente que resguarda los estándares metodológicos y los objetivos estratégicos de la Fundación:

\begin{table}[ht]
\centering
\small
\caption{Estructura de gobernanza del proyecto}
\label{tab:gobernanza}
\begin{tabular}{p{4cm}p{5.2cm}p{4.8cm}}
\hline
\textbf{Instancia} & \textbf{Integrantes} & \textbf{Rol} \\
\hline
Comité Directivo y Dirección Ejecutiva de JxR & --- & Supervisión estratégica y recepción de resultados \\
Comité AIM & --- & Revisión del proyecto \\
Panel de Revisión Metodológica y Operativa & Daniel Horta (mentor, Activa); Anastasia Ganderats (Activa); Valentina Piña (AIM y Valtia); Rodrigo de la Riva (Activa y directorio AIM) & Acompañamiento y revisión del diseño, el control y el análisis \\
Equipo de investigación & Vicente Rojas Murúa y Leonardo Capobianco (Investigadores) & Diseño, coordinación, participación en las entrevistas, procesamiento, análisis y redacción de informes \\
Equipo de profesionales & Por definir & Realización de las entrevistas \\
Comité de ética externo & Por confirmar & Evaluación ética independiente \\
\hline
\end{tabular}
\end{table}

\subsection{Roles y responsabilidades}

\begin{itemize}
    \item \textbf{Investigadores (Vicente Rojas Murúa y Leonardo Capobianco):} diseñan los instrumentos, coordinan el terreno, participan en las entrevistas y se encargan del procesamiento, el análisis y la redacción de informes.
    \item \textbf{Equipo de profesionales:} realiza las entrevistas siguiendo la pauta y los protocolos éticos, junto con los investigadores y tras una inducción previa.
    \item \textbf{Panel de revisión y Comité AIM:} revisan la metodología y el avance del proyecto.
    \item \textbf{Comité Directivo de JxR:} valida los objetivos de negocio y es la contraparte estratégica. Comparte el catastro de empresas y los contactos de mujeres, y recibe la pauta de entrevista antes del terreno.
\end{itemize}

\subsection{Gobernanza ética y resguardo de datos}

La investigación involucra a una población con alto grado de vulnerabilidad (mujeres con antecedentes penales) y a datos corporativos sensibles. Se rige por el código ICC/ESOMAR.

\begin{enumerate}
    \item \textbf{Comité de ética externo e independiente (por confirmar):} el protocolo, las pautas de entrevista y los consentimientos informados serán evaluados por un comité externo al proyecto y a JxR, sin integrantes con conflicto de interés. La revisión comienza en la Semana~1 y ninguna entrevista se realiza antes del dictamen favorable.
    \item \textbf{Consentimiento informado diferenciado:}
    \begin{itemize}
        \item \textit{Mujeres:} consentimiento explícito, gestionado directamente por el equipo investigador, que garantiza que participar, o no, no condiciona sus beneficios ni atenciones en los programas de JxR.
        \item \textit{Actores corporativos:} consentimiento informado a directivos, gerentes de RR.HH. y jefaturas, con confidencialidad de la razón social y sin vínculo con negociaciones comerciales previas o futuras con JxR.
    \end{itemize}
    \item \textbf{Protocolo para las díadas:} lo que diga la trabajadora no llega a la empresa ni al revés. Se prevén además contención para las participantes (con derivación si es necesario) y una compensación por su tiempo.
    \item \textbf{Anonimización y opt-in separado:} los datos se procesan codificados y anonimizados. La opción de contacto con JxR (opt-in) se gestiona con un formulario independiente de los instrumentos de investigación, al cierre de cada contacto. La investigación no se utiliza como prospección comercial.
\end{enumerate}
